% =====================================================================
\clearpage
\phantomsection
\section*{APÊNDICE A --- REPRODUTIBILIDADE}
\label{ap:reprodutibilidade}
\addcontentsline{toc}{section}{APÊNDICE A -- REPRODUTIBILIDADE}

\subsection*{A.1 Preparação do ambiente}

O binário do verificador acompanha o repositório, de modo que não há etapa obrigatória de
compilação para reproduzir os resultados reportados na Seção~\ref{sec:resultados}.

\begin{lstlisting}[language=bash]
git clone --recurse-submodules <url-do-repositorio>
cd esbmc/pibic
python3 -m venv .venv && source .venv/bin/activate
python -m pip install -e ".[test]"
pytest tests/ -v
\end{lstlisting}

O invólucro localiza o binário nesta ordem: a variável \texttt{\$ESBMC\_BIN}; \texttt{esbmc}
no \texttt{\$PATH}; uma compilação local em \texttt{../build/src/esbmc/esbmc}; e, por fim, o
binário versionado. Se nenhum existir, a suíte é \textbf{pulada} com motivo explícito, em vez
de acumular erros que seriam confundidos com falhas de propriedade.

\subsection*{A.2 Reexecução das evidências}

\begin{lstlisting}[language=bash]
python scripts/gerar_evidencias.py       # reexecuta os 12 harnesses
# saida bruta -> results/logs/*.log
# tabela      -> results/EVIDENCIAS.md
\end{lstlisting}

\subsection*{A.3 Verificação ilimitada por IC3/PDR}

\begin{lstlisting}[language=bash]
sudo apt-get install -y yosys            # traz o ABC embutido
cd ic3
python3 gen_transition_system.py --bits 16 -o cl_ddpg16.v
python3 validate_forward.py              # teste diferencial em 12 estados
./run_pdr.sh cl_ddpg16.v 1800
\end{lstlisting}

A execução de \texttt{validate\_forward.py} não é opcional: ela compara o Verilog gerado com
a referência em ponto fixo antes de qualquer medição. A coincidência das doze amostras detecta
regressão de tradução; não constitui prova de equivalência universal.

\subsection*{A.4 Compilação deste relatório}

\begin{lstlisting}[language=bash]
cd pibic/relatorio_final
make                 # pdflatex -> bibtex -> pdflatex -> pdflatex
make clean           # remove os intermediarios
\end{lstlisting}

As figuras são localizadas por \texttt{\textbackslash graphicspath} em \texttt{figuras/} e,
subsidiariamente, em \texttt{../artigo/figs/}. São necessários os pacotes
\texttt{texlive-latex-base}, \texttt{texlive-latex-recommended},
\texttt{texlive-latex-extra}, \texttt{texlive-fonts-recommended},
\texttt{texlive-lang-portuguese} e \texttt{texlive-publishers}, este último por fornecer
\texttt{abntex2cite} e o estilo bibliográfico \texttt{abntex2-alf}.

\subsection*{A.5 Advertências operacionais}

\begin{itemize}[leftmargin=1.4cm]
  \item \textbf{Nunca testar \texttt{not r.is\_safe} para concluir que um defeito foi
  encontrado.} Esse valor também é falso quando o verificador não executou. Use
  \texttt{r.status == UNSAFE}.
  \item \textbf{\texttt{--bounds-check} e \texttt{--pointer-check} não existem} no ESBMC:
  essas verificações são ativadas por padrão e a ferramenta expõe apenas as formas negativas,
  \texttt{--no-bounds-check} e \texttt{--no-pointer-check}. Passar as formas inexistentes
  encerra o processo com código 64.
  \item \textbf{\texttt{TIMEOUT} é indeciso}, nunca ``seguro''.
  \item \textbf{\texttt{--no-unwinding-assertions} só é seguro em \textit{harnesses} sem
  laços.} Em código com laços, o uso da opção pode tornar o caso vácuo.
\end{itemize}
